% Source-derived chapter 01: Introduction
% Original source: universal-secant-bundles-syzygies-canonical-curves
\chapter{Introduction}
\label{universal-secant-bundles-syzygies-canonical-curves-chapter-01}
\source{universal-secant-bundles-syzygies-canonical-curves}

\begin{remark}[Source section]
\label{universal-secant-bundles-syzygies-canonical-curves-chapter-01-source}
\source{universal-secant-bundles-syzygies-canonical-curves}
\source{universal-secant-bundles-syzygies-canonical-curves}
This chapter reproduces the corresponding section of the retrieved
LaTeX source, including its definitions, statements, examples, and
proofs. It has not yet been translated into Lean.
\notready
\end{remark}

% The source section title was: Introduction
In this paper, we introduce a \emph{universal} version of the secant sheaf construction from \cite{ein-lazarsfeld-asymptotic}. Using this tool, we give a simpler proof of a theorem of Voisin \cite{V1}, \cite{V2} on the equations of canonical curves. We further give a generalization of her result for curves of even genus.\smallskip

 The classical Theorem of Noether--Babbage--Petri states that canonical curves are projectively normal, and that the ideal $I_{C/\PP^{g-1}}$ is generated by quadrics (with a few exceptions), see \cite{arbarello-sernesi-petri} for a modern treatment. In the 1980s, M.\ Green realized that these classical results about the equations defining canonical curves should be the first case of a much more general statement about higher syzygies, and he made a very influential conjecture \cite{green-koszul} in this direction.\smallskip

Whilst the general case of Green's Conjecture remains open, in 2002 Voisin made a breakthrough by proving the conjecture for \emph{general} curves of even genus \cite{V1}.  Voisin's argument relies on an intricate study of the geometry of Hilbert schemes on a K3 surface. Recently, an algebraic approach to Voisin's Theorem has been given, \cite{AFPRW}, based on degenerating to the tangent developable, a singular surface whose hyperplane sections are cuspidal curves. The authors apply the representation theory of an $SL_2$ action present in this special situation to establish Green's conjecture for rational cuspidal curves. Explicit plethysm formulae play the key role, involving a change of basis between elementary symmetric polynomials and Schur polynomials. Maps which are simple to describe in one basis become rather complicated in the other, making the proof quite technical, see \cite[\S 5.5--5.7]{AFPRW}. \smallskip

%Maps which are simple to describe in one basis become rather complicated in the other, making the proof quite technical, see \cite[\S 5.5--5.7]{AFPRW}. \smallskip
%One key input for both of these proofs is an in-depth study of the geometry of the Grassmannian of lines. This has been widely considered the most ``essential" part of the %proofs. \smallskip

In this paper, we first give a simpler proof of Voisin's Theorem, using only basic homological algebra and without the need to degenerate. We further provide a structure theorem in the even genus case, describing in detail the extremal syzygy space. Let $X$ be a complex K3 surface with Picard group generated by an ample line bundle $L$ of even genus $g=2k$, i.e.\ $(L)^2=2g-2$. We proceed by direct computation on $X$. Define $\mathrm{K}_{p,q}(X,L)$ as the middle cohomology of 
{\small{$$\bigwedge^{p+1} \mathrm{H}^0(X,L) \otimes \mathrm{H}^0(X,L^{\otimes q-1}) \to \bigwedge^p \mathrm{H}^0(X,L) \otimes \mathrm{H}^0(X,L^{\otimes q}) \to \bigwedge^{p-1} \mathrm{H}^0(X,L) \otimes \mathrm{H}^0(X,L^{\otimes q+1}) $$}}
 Voisin's Theorem states that $\mathrm{K}_{k,1}(X,L)=0$, \cite{V1}. This single vanishing suffices to prove Green's Conjecture for general canonical curves in even genus. \smallskip


Let $E$ be the rank two \emph{Lazarsfeld--Mukai bundle} associated to a $g^1_{k+1}$ on a smooth curve $C \in |L|$, see \cite{lazarsfeld-BNP}. The dual bundle $E^{\vee}$ fits into the exact sequence
$$0 \to E^{\vee} \to \mathrm{H}^0(C,A) \otimes \mathcal{O}_X \to i_*A \to 0,$$
for $A$ a $g^1_{k+1}$ on $C$, where $i: C \hookrightarrow X$ is the inclusion. The vector bundle $E$ has invariants $\det(E)=L$, $h^0(E)=k+2$, $h^1(E)=h^2(E)=0$.\smallskip


We deduce Voisin's Theorem from the K\"unneth formula on $X \times \PP(\mathrm{H}^0(E))$. Our proof reduces to showing that a certain square matrix is nonsingular. Since our matrix takes the form $\mathrm{H}^k(\mathrm{Sym}^{k+1} \mathcal{G}) \to \mathrm{H}^k(\mathrm{Sym}^{k} \mathcal{G} \otimes \mathcal{G})$, for some bundle $\mathcal{G}$, the desired nonsingularity is \emph{automatic}, see the proof of Proposition \ref{auto-injectivity}.\smallskip

The starting point of Voisin's proof \cite{V1} is her \emph{Hilbert Scheme Description} of Koszul cohomology via the cohomology of tautological sheaves on the Hilbert scheme, see \cite[\S 5]{aprodu-nagel}. This description has also been used to effect in \cite{ein-lazarsfeld-gonality}. Our approach uses instead the, considerably simpler and more general, \emph{Kernel Bundle} approach to syzygies, \cite{lazarsfeld-VBT}, as our starting point.\smallskip

The projective space $\PP:=\PP(\mathrm{H}^0(E))$ can be embedded into the Hilbert scheme $X^{[k+1]}$ by sending $s\in \mathrm{H}^0(E)$ to the zero scheme $Z(s)$. Whilst $\PP \seq X^{[k+1]}$ is also a crucial object in Voisin's proof, in order to make her argument work $\PP$ needs to be replaced by the space of zero cycles of the form $Z(s)-x+y$ with $x \in Z(s)$ and $y \in X$. This results in significant complications in Voisin's argument, see \cite[\S 6.3]{aprodu-nagel}. By contrast, we find a way to work directly on $\PP$.\smallskip

  A few years after her breakthrough for even genus curves, Voisin deduced the odd genus case of generic Green's conjecture out of the even genus case \cite{V2}. In Section \ref{odd-genus}, we give a streamlined version of Voisin's argument:
  \begin{thm}
\source{universal-secant-bundles-syzygies-canonical-curves}
\notready \label{generic-green}
\source{universal-secant-bundles-syzygies-canonical-curves}
  Let $C$ be a general curve over $\C$ of genus $g=2k$ or $g=2k+1$. Then $\mathrm{K}_{k,1}(C,\omega_C)=0$.
  \end{thm}
\smallskip

 Whilst we use several of the same objects from \cite{V2}, such as nodal K3 surfaces, our proof ends us being more economical. In particular, we do away with the explicit computations of \cite[\S 2]{V2} as well as the difficult analysis of the geometry of the Grassmannian from \cite[\S 3, ``Fourth Step'']{V2}, considered by Voisin to be the most essential part of her proof. Furthermore, the application of our techniques to nodal K3 surfaces of even genus is far simpler in our setting.\smallskip

In fact, our approach proves much more than Voisin's theorem for even genus curves. Let $(X,L)$ be any polarized variety. A natural question, with roots going back to Andreotti--Mayer's 1967 paper \cite{andreotti-mayer}, is whether one can find a spanning set of the spaces $\mathrm{K}_{p,1}(X,L)$ consisting of elements of low \emph{rank}. Here one defines the rank of a syzygy $\alpha \in \mathrm{K}_{p,1}(X,L)$ as the dimension of the minimal subspace
$V \seq \mathrm{H}^0(X,L)$ such that $\alpha \in \mathrm{K}_{p,1}(X,L,V)$, i.e.\ such that $\alpha$ is represented by an element of $\bigwedge^p V \otimes \mathrm{H}^0(X,L)$, \cite{bothmer-JPAA}. Syzygies of low rank have geometric meaning. For $\alpha \in \mathrm{K}_{p,1}(X,L)$ we always have $\rm{rank}(\alpha) \geq p+1$. If there exists a syzygy $\alpha$ with $\rm{rank}(\alpha)=p+1$, then $X$ lies on a rational normal scroll. More precisely, the \emph{syzygy scheme} $$\rm{Syz}(\alpha) \seq \PP(\mathrm{H}^0(L)^{\vee})$$ of $\alpha$, as defined by Green \cite{green-canonical}, defines a scroll. Similarly, syzygies of rank $p+2$ arise from linear sections of Grassmannians. Precisely, $\rm{Syz}(\alpha)$ contains the cone over a section of $\rm{Gr}_2(V^{\vee} \oplus \C)$. \smallskip

In Section \ref{sec2}, we prove the following structure theorem for the last nonzero syzygy space for K3 surfaces of even genus.
\begin{thm}
\source{universal-secant-bundles-syzygies-canonical-curves}
\notready \label{geo-thm-intro}
\source{universal-secant-bundles-syzygies-canonical-curves}
Let $X$ be a complex K3 surface. Assume $\mathrm{Pic}(X)=\mathbb{Z}[L]$ with $L$ ample of even genus $g=2k$. Let $E$ be the Lazarsfeld-Mukai bundle as above. For any nonzero $s \in \mathrm{H}^0(E)$, the space $\mathrm{K}_{k-1,1}(X,L,\mathrm{H}^0(L \otimes I_{Z(s)}))$ is a one-dimensional subspace of $\mathrm{K}_{k-1,1}(X,L)$.
The morphism
{\small{\begin{align*}
\psi \; : \; \PP(\mathrm{H}^0(E)) & \to \PP(\mathrm{K}_{k-1,1}(X,L))\\
[s] & \mapsto [\mathrm{K}_{k-1,1}(X,L,\mathrm{H}^0(L \otimes I_{Z(s)}))]
\end{align*}}}
is the Veronese embedding of degree $k-2$. In particular, $\psi$ induces a natural isomorphism $\mathrm{Sym}^{k-2}\mathrm{H}^0(X,E) \simeq \mathrm{K}_{k-1,1}(X,L) $.
\end{thm}
A consequence of Theorem \ref{geo-thm-intro} is that $\mathrm{K}_{k-1,1}(X,L)$ is spanned by syzygies of rank $k+1=\dim \mathrm{H}^0(L \otimes I_{Z(s)})$. Theorem \ref{geo-thm-intro} implies the $\mathrm{K}_{k,1}(X,L)=0$ by standard dimension computations, \cite[\S 4.1]{farkas-progress}, and thus enhances Voisin's Theorem in even genus. One may compare Theorem \ref{geo-thm-intro} to Schreyer's Conjecture, proven in \cite{lin-syz}, describing the structure of the last syzygy space for curves of \emph{non-maximal} gonality.\smallskip

% It is possible to describe the map $\mathrm{Sym}^{k-2}\mathrm{H}^0(X,E) \to \mathrm{K}_{k-1,1}(X,L)$ explicitly, cf.\ \cite[\S 3.4]{aprodu-nagel}, however we will not need this.\smallskip 

Theorem \ref{geo-thm-intro} implies a previously open conjecture known as the \emph{Geometric Syzygy Conjecture} in even genus, see \cite{bothmer-Transactions} where the statement is proven for $g \leq 8$.  To put this conjecture in context, recall the following important result:
 \begin{thm}[\cite{andreotti-mayer}, \cite{green-quadrics}]
\source{universal-secant-bundles-syzygies-canonical-curves}
\notready \label{green-quads}
\source{universal-secant-bundles-syzygies-canonical-curves}
 The ideal $I_{C/\PP^{g-1}}$ of a canonical curve of Clifford index at least two is generated by quadrics of rank four.
 \end{thm}
 This provides an enhancement of Petri's theorem stating that $I_{C/\PP^{g-1}}$ is generated by quadrics if $\mathrm{Cliff}(C) \geq 2$. More precisely, let $W^1_{g-1}(C)$ be the locus of line bundles $A$ of degree $g-1$ with two sections; geometrically this is the locus of double points of the Theta Divisor $\Theta \seq \mathrm{Pic}^{g-1}(C)$. Then $I_{C/\PP^{g-1}}$ is generated by the quadrics defined by the Petri map
 $$\mathrm{H}^0(A) \otimes \mathrm{H}^0(\omega_C \otimes A^{-1}) \to \mathrm{H}^0(\omega_C)$$
 for $A \in W^1_{g-1}(C)$. To extend Theorem \ref{green-quads} to higher syzygies, note that Theorem \ref{green-quads} can be rephrased as stateing that $\mathrm{K}_{1,1}(C,\omega_C)$ is spanned by the spaces $\mathrm{K}_{1,1}(C,\omega_C, \mathrm{H}^0(\omega_C \otimes A^{-1} ))$ for $A \in W^1_{g-1}(C)$, and hence by syzygies of rank \emph{two}.\smallskip


If we restrict the subspaces $\mathrm{K}_{k-1,1}(X,L,\mathrm{H}^0(L \otimes I_{Z(s)})) \seq \mathrm{K}_{k-1,1}(X,L)$ to a curve $C \in |L|$ containing the locus $Z(s)$, then $\mathrm{K}_{k-1,1}(X,L,\mathrm{H}^0(L \otimes I_{Z(s)}))$ restricts to 
$$\mathrm{K}_{k-1,1}(C,\omega_C,\mathrm{H}^0(\omega_C \otimes A^{-1})) \seq \mathrm{K}_{k-1,1}(C,\omega_C)$$ where $A \in W^1_{k+1}(C)$ is the line bundle defined by the divisor $Z(s) \seq C$. Thus, the rank $k+1$ syzygies $\alpha \in \mathrm{K}_{k-1,1}(X,L,\mathrm{H}^0(L \otimes I_{Z(s)}))$ drop rank by one when restricted to $C$. Combining this with an argument of Voisin,  see \cite[Prop.\ 7]{V1} and the unpublished \cite[\S 11]{bothmer-preprint} we obtain:
\begin{cor}[Geometric Syzygy Conjecture in Even Genus]
\source{universal-secant-bundles-syzygies-canonical-curves}
\notready \label{geo-cor}
\source{universal-secant-bundles-syzygies-canonical-curves}
Let $C$ be a general curve of even genus $g=2k$. Then $\mathrm{K}_{k-1,1}(C,\omega_C)$ is generated by syzygies of the lowest possible rank $k$.  More precisely, $\mathrm{K}_{k-1,1}(C,\omega_C)$ is generated by the rank $k$ syzygies $$\alpha \in \mathrm{K}_{k-1,1}(C,\omega_C,\mathrm{H}^0(\omega_C \otimes A^{-1})), \; \; A \in W^1_{k+1}(C).$$
\end{cor} 
Corollary \ref{geo-cor} therefore provides an extension of Green's theorem on quadrics \cite{green-quadrics} to the space of linear syzygies of highest order.\smallskip 

It would appear to us to be very difficult to adapt degeneration methods to prove the structure Theorem \ref{geo-thm-intro}, as opposed to merely establishing the vanishing from Voisin's original result. For instance, the construction of Lazarsfeld--Mukai bundles fails on the (non-normal) tangent developable, so that it is not even clear how the bundle $E$ degenerates to this surface. \smallskip

There are no known conjectural candidates for an analogous result to Theorem \ref{geo-thm-intro} in odd genus $g=2k+1$. In this case, the dimension of $\mathrm{K}_{k-1,1}(X,L)$ is not given by a binomial coefficient, so this space cannot be of the form $\mathrm{Sym}^p(V)$ for any vector space $V$. Furthermore, we no longer have uniqueness of the relevant Lazarsfeld-Mukai bundle in this situation.\smallskip

Our argument is formal, using general results on vector bundles rather than a detailed study of the geometry of curves. Consequently, we expect our approach to generalize well to higher dimensional varieties, for which previous approaches to Green's conjecture seem less applicable. See \cite{ein-laz-arb-dim} for applications of vector bundle methods to syzygies of varieties of high dimension.\smallskip
 
 %There do not exist K3 surfaces of Picard rank one over the algebraic closure of a finite field, \cite[Ch.\ 17]{huybrechts-k3}. However, as was explained to us by J.\ Rathmann, we may still use Theorem \ref{geo-thm-intro} to deduce the Voisin's Theorem \ref{generic-green} for a general curve of even genus $g=2k$ over \emph{any} algebraically closed field $\mathbb{F}$ of characteristic $p \geq \frac{g+4}{2}$. Indeed, 
 %the geometric general fibre $(X_{\bar{\tau}}, L_{\bar{\tau}})$ of the universal deformation of any polarized K3 surface $(X, L)$ over $k$ has Picard rank one, \cite[Thm.\ 2.9]{ogus}. Theorem \ref{geo-thm-intro} applies to the geometric general fibre $X_{\bar{\tau}}$; in particular    $K_{k,1}(X_{\bar{\tau}},L_{\bar{\tau}})=0$. Semicontinuity of Koszul cohomology, together with flat base change, implies that there exist polarized K3 surfaces $(X',L')$ over the \emph{original} field $\mathbb{F}$ with $K_{k,1}(X',L')=0$, as required. \smallskip
 
 
% Additional arguments are needed in order to extend our proof to cover the odd genus case of Voisin's Theorem in positive characteristic. For instance, one would need an argument for the existence of the K3 surface of Picard rank two from Section \ref{odd-genus}.\smallskip

%Green's Conjecture applied to general curves is particularly strong in that it completely determines the terms in the resolution of the canonical ring $\Gamma_C(\omega_C):=\bigoplus_q \mathrm{H}^0(\omega_C^{\otimes q})$, see \cite[\S 4.1]{farkas-progress}. It is very likely that our proof works in characteristic $p$ for large $p$, but we have not checked this.\smallskip


 %The reader may notice a resemblance between our argument and \cite{V2}, Proof of Proposition 8. The idea of considering $\text{Sym}^{k+1} \mathcal{S}$ was %foreshadowed in \cite{AFPRW}. Other contributions to this topic include \cite{aprodu-farkas}, which relies on Voisin's results, and \cite{raicu-sam}, which relies upon %\cite{AFPRW}.\medskip



\textbf{Acknowledgements} I thank C.\ Voisin for helpful explanations and G.\ Farkas for numerous discussions. I thank R.\ Lazarsfeld for encouragement and for detailed comments. I thank M.\ Aprodu, J.\ Ellenberg, D.\ Erman, D.\ Huybrechts, J.\ Rathmann, E.\ Sernesi and R.\ Yang for feedback on previous versions.\smallskip

I thank the referee for a careful reading and for comments which greatly improved the exposition. The author is supported by NSF grant DMS-1701245.

\subsection{Preliminaries}
We work over the complex numbers throughout. We gather here a few facts. We have a natural isomorphism $(\mathrm{Sym}^j F)^{\vee} \simeq \mathrm{Sym}^j F^{\vee}$ for any vector bundle $F$ on a variety $X$ defined over $\C$ (over arbitrary fields, this statement requires the characteristic to be at least $j+1$). Let $0 \to F_1 \to F_2 \to F_3 \to 0$ be an exact sequence of bundles on $X$. If $p=0$ or $p \geq i+1$ then, 
from \cite[\S V]{ABW}, we have exact sequences
{\small{\begin{align*}
&\ldots \to \bigwedge^{i-2} F_2 \otimes \mathrm{Sym}^2(F_1) \to \bigwedge^{i-1} F_2 \otimes F_1  \to \bigwedge^i F_2 \to \bigwedge^i F_3 \to 0, \\
& \ldots \to \mathrm{Sym}^{i-2}F_2 \otimes \bigwedge^2 F_1 \to \mathrm{Sym}^{i-1}F_2 \otimes F_1 \to \mathrm{Sym}^{i}F_2 \to \mathrm{Sym}^{i}F_3 \to 0. 
\end{align*}}}
 We may further dualize $0 \to F_1 \to F_2 \to F_3 \to 0$, form the second exact sequence above and dualize again to obtain an exact sequence
{\small{$$
		0 \to \mathrm{Sym}^i(F_1) \to \mathrm{Sym}^{i}(F_2) \to \mathrm{Sym}^{i-1}(F_2) \otimes F_3 \to \mathrm{Sym}^{i-2}(F_2) \otimes \bigwedge^2 F_3 \to \ldots
$$}}
Let $f: X \to Y$ be a morphism of varieties and $\mathcal{F} \in \mathrm{Coh}(X)$ a sheaf. If $\mathcal{E}$ is a vector bundle on $Y$ then we have the \emph{Projection Formula} $\mathrm{R}^if_*(\mathcal{F} \otimes f^* \mathcal{E}) \simeq \mathrm{R}^if_*\mathcal{F} \otimes \mathcal{E}$, \cite[III, Ex.\ 8.3]{hartshorne}. In particular, $f_*(\mathcal{F} \otimes f^* \mathcal{E}) \simeq f_*\mathcal{F} \otimes \mathcal{E}$. If $\mathrm{R}^if_*\mathcal{F}=0$ for all $ i>0$ then $\mathrm{H}^p(X,\mathcal{F}) \simeq \mathrm{H}^p(Y, f_* \mathcal{F})$ for $p \geq 0$, \cite[III, Ex.\ 8.1 ]{hartshorne}. If $X, Y$ are varieties and $\mathcal{F} \in \mathrm{Coh}(X), \mathcal{G} \in \mathrm{Coh}(Y)$ are sheaves, the \emph{K\"unneth formula} states
 {\small{$$\mathrm{H}^{m}(X \times Y, \mathcal{F} \boxtimes \mathcal{G}) \simeq \bigoplus_{a+b=m} \mathrm{H}^a(X,\mathcal{F}) \otimes \mathrm{H}^b(Y, \mathcal{G}),$$}}
where $\mathcal{F} \boxtimes \mathcal{G}:=p^*\mathcal{F} \otimes q^* \mathcal{G}$, for projections $p: X \times Y \to X$, $q: X \times Y \to Y$.\smallskip
 
 Assume we have an exact sequence $0 \to \mathcal{F} \to \mathcal{G} \to \mathcal{H} \to 0$ of coherent sheaves on a quasi-projective variety, with $\mathcal{G}$ locally free. Assume either $\mathcal{H}$ is locally free or $\mathcal{H} \simeq \mathcal{O}_D$ for a Cartier divisor $D$. Then $\mathcal{F}$ is locally free. This follows from \cite[III, Ex 6.5]{hartshorne}.
